% Options for packages loaded elsewhere
\PassOptionsToPackage{unicode}{hyperref}
\PassOptionsToPackage{hyphens}{url}
\PassOptionsToPackage{dvipsnames,svgnames,x11names}{xcolor}
%
\documentclass[
  10pt,
]{article}
\usepackage{amsmath,amssymb}
\usepackage{iftex}
\ifPDFTeX
  \usepackage[T1]{fontenc}
  \usepackage[utf8]{inputenc}
  \usepackage{textcomp} % provide euro and other symbols
\else % if luatex or xetex
  \usepackage{unicode-math} % this also loads fontspec
  \defaultfontfeatures{Scale=MatchLowercase}
  \defaultfontfeatures[\rmfamily]{Ligatures=TeX,Scale=1}
\fi
\usepackage{lmodern}
\ifPDFTeX\else
  % xetex/luatex font selection
  \setmainfont[]{DejaVu Serif}
  \setsansfont[]{DejaVu Sans}
  \setmonofont[]{DejaVu Sans Mono}
  \setmathfont[]{Latin Modern Math}
\fi
% Use upquote if available, for straight quotes in verbatim environments
\IfFileExists{upquote.sty}{\usepackage{upquote}}{}
\IfFileExists{microtype.sty}{% use microtype if available
  \usepackage[]{microtype}
  \UseMicrotypeSet[protrusion]{basicmath} % disable protrusion for tt fonts
}{}
\makeatletter
\@ifundefined{KOMAClassName}{% if non-KOMA class
  \IfFileExists{parskip.sty}{%
    \usepackage{parskip}
  }{% else
    \setlength{\parindent}{0pt}
    \setlength{\parskip}{6pt plus 2pt minus 1pt}}
}{% if KOMA class
  \KOMAoptions{parskip=half}}
\makeatother
\usepackage{xcolor}
\usepackage[margin=2.0cm]{geometry}
\usepackage{color}
\usepackage{fancyvrb}
\newcommand{\VerbBar}{|}
\newcommand{\VERB}{\Verb[commandchars=\\\{\}]}
\DefineVerbatimEnvironment{Highlighting}{Verbatim}{commandchars=\\\{\}}
% Add ',fontsize=\small' for more characters per line
\newenvironment{Shaded}{}{}
\newcommand{\AlertTok}[1]{\textcolor[rgb]{1.00,0.00,0.00}{\textbf{#1}}}
\newcommand{\AnnotationTok}[1]{\textcolor[rgb]{0.38,0.63,0.69}{\textbf{\textit{#1}}}}
\newcommand{\AttributeTok}[1]{\textcolor[rgb]{0.49,0.56,0.16}{#1}}
\newcommand{\BaseNTok}[1]{\textcolor[rgb]{0.25,0.63,0.44}{#1}}
\newcommand{\BuiltInTok}[1]{\textcolor[rgb]{0.00,0.50,0.00}{#1}}
\newcommand{\CharTok}[1]{\textcolor[rgb]{0.25,0.44,0.63}{#1}}
\newcommand{\CommentTok}[1]{\textcolor[rgb]{0.38,0.63,0.69}{\textit{#1}}}
\newcommand{\CommentVarTok}[1]{\textcolor[rgb]{0.38,0.63,0.69}{\textbf{\textit{#1}}}}
\newcommand{\ConstantTok}[1]{\textcolor[rgb]{0.53,0.00,0.00}{#1}}
\newcommand{\ControlFlowTok}[1]{\textcolor[rgb]{0.00,0.44,0.13}{\textbf{#1}}}
\newcommand{\DataTypeTok}[1]{\textcolor[rgb]{0.56,0.13,0.00}{#1}}
\newcommand{\DecValTok}[1]{\textcolor[rgb]{0.25,0.63,0.44}{#1}}
\newcommand{\DocumentationTok}[1]{\textcolor[rgb]{0.73,0.13,0.13}{\textit{#1}}}
\newcommand{\ErrorTok}[1]{\textcolor[rgb]{1.00,0.00,0.00}{\textbf{#1}}}
\newcommand{\ExtensionTok}[1]{#1}
\newcommand{\FloatTok}[1]{\textcolor[rgb]{0.25,0.63,0.44}{#1}}
\newcommand{\FunctionTok}[1]{\textcolor[rgb]{0.02,0.16,0.49}{#1}}
\newcommand{\ImportTok}[1]{\textcolor[rgb]{0.00,0.50,0.00}{\textbf{#1}}}
\newcommand{\InformationTok}[1]{\textcolor[rgb]{0.38,0.63,0.69}{\textbf{\textit{#1}}}}
\newcommand{\KeywordTok}[1]{\textcolor[rgb]{0.00,0.44,0.13}{\textbf{#1}}}
\newcommand{\NormalTok}[1]{#1}
\newcommand{\OperatorTok}[1]{\textcolor[rgb]{0.40,0.40,0.40}{#1}}
\newcommand{\OtherTok}[1]{\textcolor[rgb]{0.00,0.44,0.13}{#1}}
\newcommand{\PreprocessorTok}[1]{\textcolor[rgb]{0.74,0.48,0.00}{#1}}
\newcommand{\RegionMarkerTok}[1]{#1}
\newcommand{\SpecialCharTok}[1]{\textcolor[rgb]{0.25,0.44,0.63}{#1}}
\newcommand{\SpecialStringTok}[1]{\textcolor[rgb]{0.73,0.40,0.53}{#1}}
\newcommand{\StringTok}[1]{\textcolor[rgb]{0.25,0.44,0.63}{#1}}
\newcommand{\VariableTok}[1]{\textcolor[rgb]{0.10,0.09,0.49}{#1}}
\newcommand{\VerbatimStringTok}[1]{\textcolor[rgb]{0.25,0.44,0.63}{#1}}
\newcommand{\WarningTok}[1]{\textcolor[rgb]{0.38,0.63,0.69}{\textbf{\textit{#1}}}}
\usepackage{longtable,booktabs,array}
\usepackage{calc} % for calculating minipage widths
% Correct order of tables after \paragraph or \subparagraph
\usepackage{etoolbox}
\makeatletter
\patchcmd\longtable{\par}{\if@noskipsec\mbox{}\fi\par}{}{}
\makeatother
% Allow footnotes in longtable head/foot
\IfFileExists{footnotehyper.sty}{\usepackage{footnotehyper}}{\usepackage{footnote}}
\makesavenoteenv{longtable}
\setlength{\emergencystretch}{3em} % prevent overfull lines
\providecommand{\tightlist}{%
  \setlength{\itemsep}{0pt}\setlength{\parskip}{0pt}}
\setcounter{secnumdepth}{-\maxdimen} % remove section numbering
% doc_header.tex -- LaTeX preamble additions used by build_docs.sh for every
% Markdown -> LaTeX -> PDF conversion (the documents of the folder Revision_HS_CJT_and_Kohn-Sham).
\usepackage{etoolbox}
% longtable may not be loaded by pandoc when a document has no table
\usepackage{longtable}
% bold vectors: \bm -> \symbf under unicode-math (XeLaTeX), bm package otherwise
\makeatletter
\@ifpackageloaded{unicode-math}{\providecommand{\bm}[1]{\symbf{#1}}}{\usepackage{bm}}
\makeatother
% smaller type in long tables so that wide audit tables fit the text width;
% inside a table a line break is allowed after an underscore of plain text
% (pandoc writes it as \_), so that long identifiers wrap inside their column
% rows are set a little apart, so that symbols with indices do not touch
\AtBeginEnvironment{longtable}{\footnotesize\renewcommand{\arraystretch}{1.45}%
  \let\sptunderscore\_\def\_{\sptunderscore\allowbreak}}
\setlength{\LTcapwidth}{\textwidth}
% A short table is kept on one page.  The build (code/scripts/keep_tables.py) ends the rows of such a table with \\*,
% which forbids a page break after a row, and sets \sptkeptpenalty to 10000 in front of the table.  This penalty is
% put here in the place of the breakpoint (penalty 0) that longtable leaves after every table, so that the paragraph
% that follows a kept table, its caption, begins on the same page; it is set back to 0 after the table.
\newcount\sptkeptpenalty
\sptkeptpenalty=0
\makeatletter
\patchcmd{\endlongtable}{\endgraf\penalty\z@\addvspace\LTpost}%
  {\endgraf\penalty\sptkeptpenalty\addvspace\LTpost\global\sptkeptpenalty=\z@}%
  {}{\PackageError{doc_header}{the end of longtable could not be patched}{The tables can still be typeset, but a short table may be separated from its caption.}}
\makeatother
% allow line breaks after slashes and hyphens inside table cells
\usepackage{xurl}
% code blocks: break over-long lines instead of running into the margin.
% Plain ``` blocks are typeset by pandoc as {verbatim}, highlighted ones
% (```bash) as {Highlighting}; both are re-declared as breakable fancyvrb
% environments.  The continuation line is marked with a small hooked arrow.
\usepackage{fvextra}
\RecustomVerbatimEnvironment{verbatim}{Verbatim}{breaklines,breakanywhere,fontsize=\small}
\makeatletter
\@ifundefined{Highlighting}{}{%
  \RecustomVerbatimEnvironment{Highlighting}{Verbatim}{breaklines,breakanywhere,fontsize=\small,commandchars=\\\{\}}}
\makeatother
% inline code (pandoc writes it as \texttt{...}): never hyphenate -- an inserted
% hyphen could be read as part of a file name -- but allow a line break after
% an underscore, slash, comma, colon or equals sign, so that long paths and
% column lists wrap inside the text width and inside table columns.
% Inside a formula (e.g. the name of a notebook symbol in an equation) the same
% command must first leave math mode, which \text of amsmath does.
\ExplSyntaxOn
\tl_new:N \l__spt_code_tl
\cs_new_protected:Npn \__spt_texttt:n #1
  {
    \tl_set:Nn \l__spt_code_tl {#1}
    \regex_replace_all:nnN { ([\/\,\:\=]) } { \1 \c{allowbreak} } \l__spt_code_tl
    \regex_replace_all:nnN { \c{_} } { \c{_} \c{allowbreak} } \l__spt_code_tl
    \group_begin: \ttfamily \hyphenchar \font = -1 \scan_stop: \l__spt_code_tl \group_end:
  }
\RenewDocumentCommand \texttt { m }
  {
    \mode_if_math:TF { \text { \__spt_texttt:n {#1} } } { \__spt_texttt:n {#1} }
  }
\ExplSyntaxOff
% A figure is written in the Markdown source as an image paragraph followed by
% a caption paragraph ("**Figure S1.** ..."), so that the caption is also
% visible on GitHub.  An image that starts a paragraph is centered.
\usepackage{graphicx}% pandoc loads it only for documents that contain an image
\NewCommandCopy\sptincludegraphics\includegraphics
\RenewDocumentCommand\includegraphics{O{} m}{%
  \ifvmode\begin{center}\sptincludegraphics[#1]{#2}\end{center}%
  \else\sptincludegraphics[#1]{#2}\fi}
% keep figures where they are referenced
\usepackage{float}
\floatplacement{figure}{H}
% running header with the document title
\usepackage{fancyhdr}
\pagestyle{fancy}
\fancyhf{}
\fancyhead[L]{\small PROJECT spacetimes, Revision\_HS\_CJT\_and\_Kohn-Sham}
\fancyhead[R]{\small\thepage}
\renewcommand{\headrulewidth}{0.4pt}
% 16x16 matrices: tighter columns so that a full matrix fits the text width
\setlength{\arraycolsep}{3.2pt}
% The documents of this project contain long tagged symbols, file paths and
% identifiers that cannot be hyphenated.  A paragraph is therefore allowed to
% have looser lines before TeX gives up and lets a line stick out of the margin.
\AtBeginDocument{\tolerance=2000 \emergencystretch=5em \hbadness=2000}
\ifLuaTeX
  \usepackage{selnolig}  % disable illegal ligatures
\fi
\IfFileExists{bookmark.sty}{\usepackage{bookmark}}{\usepackage{hyperref}}
\IfFileExists{xurl.sty}{\usepackage{xurl}}{} % add URL line breaks if available
\urlstyle{same}
\hypersetup{
  pdftitle={paper\_prd: the Regular Article ``Higgs as a Composite, and Quantum Gravity'' for Physical Review D},
  pdfauthor={Prepared for Patrick Lee Nash, Ph.D.},
  colorlinks=true,
  linkcolor={blue},
  filecolor={Maroon},
  citecolor={blue},
  urlcolor={blue},
  pdfcreator={LaTeX via pandoc}}

\title{paper\_prd: the Regular Article ``Higgs as a Composite, and
Quantum Gravity'' for Physical Review D}
\usepackage{etoolbox}
\makeatletter
\providecommand{\subtitle}[1]{% add subtitle to \maketitle
  \apptocmd{\@title}{\par {\large #1 \par}}{}{}
}
\makeatother
\subtitle{README of the folder
Revision\_HS\_CJT\_and\_Kohn-Sham/paper\_prd of PROJECT spacetimes}
\author{Prepared for Patrick Lee Nash, Ph.D.}
\date{November 21, 2026}

\begin{document}
\maketitle

\hypertarget{what-this-folder-is}{%
\section{1. What this folder is}\label{what-this-folder-is}}

This folder contains the Regular Article ``Higgs as a Composite, and
Quantum Gravity'' by Patrick Lee Nash, Ph.D., prepared for submission to
Physical Review D, together with its Supplemental Material, the cover
letter, the information to be entered on the submission server of the
American Physical Society (APS), and the folder of the files to be
uploaded. The author asked for it on 21 November 2026 at 11:29 PST, with
three questions:

\begin{enumerate}
\def\labelenumi{\arabic{enumi}.}
\tightlist
\item
  How wide is the region of the fluctuations near the critical
  temperature \(T^{X,{\rm c}}_{\rm SF}\) of the superfluid? The answer
  needs one more calculation, the gradient coefficient of the order
  parameter at \(T^{X,{\rm c}}_{\rm SF}\).
\item
  For each of the five mean-field methods \(X\) = HS \((1,1)\), HS
  \((0,1)\), HS \((0,0)\), KS, CJT: can the superfluid play the role of
  the Higgs particle?
\item
  Is the Higgs then a composite particle?
\end{enumerate}

The author also fixed the form of the answer: a Regular Article for
Physical Review D, as a publication-ready PDF with its text, references,
figures, tables and captions, its LaTeX source in REVTeX, a data
availability statement, a cover letter, generous supplementary documents
that prove every assertion, and the preparation of the steps on the
submission server. The steps on the server (logging in, entering the
data, uploading, submitting) are the author's; nothing has been
submitted. The gravitational sector of the title is not computed: the
article outlines in its Sec. VIII how it will be added through E. A.
Lord's gauge theory of the local change of spinor representation (Proc.
Cambridge Philos. Soc. 69, 423 (1971)), the route that the author named,
and says that this sector is a program.

\hypertarget{the-answers-in-brief}{%
\section{2. The answers in brief}\label{the-answers-in-brief}}

The symbols are those of the article: \(X\) is the method,
\(\mathrm{SF}\) marks a quantity of the superfluid, \(\Lambda_{\rm LL}\)
(the momentum cutoff of the uniform states) is the unit of energy,
\(v^X_{\rm SF}=2f^{X,{\rm I}}_{\rm SF}\) is the scale of the breaking of
the new gauge group, defined as the Standard Model defines its vacuum
value, by \(M_W=gv/2\) (\(f^{X,{\rm I}}_{\rm SF}\) the decay constant of
the isospin rotations; the value of \(v^X_{\rm SF}\) in GeV is not
known), and \(M^{X,{\rm H}}_{\rm SF}\) is the mass of the amplitude
mode. The spinor field is the DiracLike spinor: it has electric charge 0
and no electroweak quantum numbers, and
\(SU(2)_{\rm v}\times U(1)^{\rm F}_{\rm LL}\) is a new gauge group, not
the electroweak group.

\begin{enumerate}
\def\labelenumi{\arabic{enumi}.}
\tightlist
\item
  \textbf{The width of the critical region.} The Ginzburg number, the
  reduced temperature
  \(\lvert T-T^{X,{\rm c}}_{\rm SF}\rvert/T^{X,{\rm c}}_{\rm SF}\) below
  which the Gaussian fluctuations of the order parameter exceed the jump
  of the specific heat of the mean-field theory, is
  \({\rm Gi}^{X}_{\rm SF}\) = \(7.6\) (KS, CJT) to \(11\) (HS (1, 1)).
  The critical region is wider than the critical temperature itself, and
  the mean-field transition, on which the five methods agree, is not
  reliable. The cause is that the pairs are tightly bound: the
  correlation length is \(1.05\) (HS (1, 1)) to \(1.35\) (HS (0, 0)) in
  units of \(1/\Lambda_{\rm LL}\), about the cutoff length, and a
  correlation volume holds \(0.035\) (HS (1, 1)) to \(0.043\) (KS, CJT)
  degrees of freedom. The gradient coefficient needed for this number is
  computed from the static pair response of a moving pair; the
  calculation also shows that the cutoff must be placed on the relative
  momentum of the pair (with a cutoff on the momenta of the two
  constituents the gradient expansion does not exist).
\item
  \textbf{The Higgs particle.} (a) In every method the superfluid is the
  Higgs field of its own gauge group: \(W^\pm\) and \(Z\) become
  massive, one gauge boson, the zeugon, stays massless, and a scalar of
  charge \(Q_{\rm LL}=0\) lies below the threshold of two
  quasi-particles. That scalar is a Higgs particle only in the
  mean-field dynamics: once it can decay into the longitudinal gauge
  bosons its tree-level width is \(2.1\) (KS, CJT) to \(3.9\) (HS (1,
  1)) times its mass; \(M^{X,{\rm H}}_{\rm SF}/v^X_{\rm SF}\) = \(4.46\)
  (CJT) to \(5.69\) (HS (1, 1)), the quartic coupling is \(9.94\) (CJT)
  to \(16.2\) (HS (1, 1)), and every method has the two non-unitary
  scalars of the charges \(Q_{\rm LL}=\pm1\) at \(0.991\) (HS (1, 1)) to
  \(1.09\) (KS) times \(M^{X,{\rm H}}_{\rm SF}\). The methods differ in
  degree, not in kind.

  \begin{enumerate}
  \def\labelenumii{(\alph{enumii})}
  \setcounter{enumii}{1}
  \tightlist
  \item
    In no method can the superfluid be the Higgs field of the Standard
    Model: the DiracLike spinor carries no electroweak quantum numbers,
    so the condensate is a singlet of the electroweak group, whatever
    its numbers.
  \end{enumerate}
\item
  \textbf{Compositeness.} Yes, in every method and by construction: the
  field of the superfluid is the amplitude of the pair self-energy, a
  composite of two spinor charges; its kinetic term consists of
  susceptibilities of the spinor field; its amplitude mode is a bound
  state of two quasi-particles. Whether the Higgs boson of nature is
  composite is not decided by this computation, because the Higgs
  particle of this model is not that boson.
\item
  \textbf{The new massless gauge boson.} After the breaking exactly one
  gauge boson stays massless, the gauge boson of the unbroken isospin
  charge \(Q_{\rm LL}\); it is not the photon, and the author offered
  its naming: it is the \textbf{zeugon} (Greek \emph{zeugos}, a yoked
  pair), the only gauge boson that attracts the two members of a pair.
  The surviving group \(O(2)\) makes its charge conjugation a gauge
  symmetry, the structure of Alice electrodynamics; the Alice strings
  and monopoles that follow from the topology of the vacuum are not
  computed.
\end{enumerate}

\hypertarget{contents}{%
\section{3. Contents}\label{contents}}

\begin{longtable}[]{@{}
  >{\raggedright\arraybackslash}p{(\columnwidth - 2\tabcolsep) * \real{0.5000}}
  >{\raggedright\arraybackslash}p{(\columnwidth - 2\tabcolsep) * \real{0.5000}}@{}}
\toprule\noalign{}
\begin{minipage}[b]{\linewidth}\raggedright
file or folder
\end{minipage} & \begin{minipage}[b]{\linewidth}\raggedright
what it is
\end{minipage} \\
\midrule\noalign{}
\endhead
\bottomrule\noalign{}
\endlastfoot
\texttt{higgs\_composite\_qg.pdf}, \texttt{.tex}, \texttt{.md} & the
article (REVTeX 4.2, two columns), its LaTeX source and its Markdown
version \\
\texttt{higgs\_composite\_qg\_SM.pdf}, \texttt{.tex}, \texttt{.md} & the
Supplemental Material: notation, model, methods, the superfluid at
\(T=0\) and at finite temperature, the derivation of the gradient
coefficient and of the Ginzburg number, the Higgs sector, compositeness,
Lord's construction, every check of the three implementations with its
deviation and tolerance, the data files and programs with their SHA-256
checksums \\
\texttt{cover\_letter.pdf}, \texttt{.tex}, \texttt{.md} & the cover
letter to the editors \\
\texttt{submission\_info.pdf}, \texttt{.tex}, \texttt{.md} & what to
enter on the submission server: journal and article type, length, title
and byline, the abstract (typeset and as LaTeX source), the subject
classification (PhySH), the data availability statement, the comments
for the editor and the suggested referees, the files to upload, the
steps on the server, and the statements that only the author can
confirm \\
\texttt{submission/} & exactly the files to upload: the article as PDF
and as LaTeX source with its \texttt{.bbl} and its five figures in the
same folder, the Supplemental Material and the cover letter \\
\texttt{src/} & the templates of the four documents
(\texttt{*.template.tex}) and the bibliography \texttt{references.bib};
every number that a program computes is a placeholder
\texttt{\{\{NAME\}\}} in the templates \\
\texttt{scripts/build\_paper.py} & the program that fills the
placeholders from the data files and the reports of the checks and
builds the four documents and the folder \texttt{submission/} \\
\texttt{figures/} & the five figures of the article, copied by the build
from the folder \texttt{figures/} of
\texttt{Revision\_HS\_CJT\_and\_Kohn-Sham} \\
\texttt{generated/numbers.csv} & every placeholder with the value that
the last build inserted \\
\texttt{generated/build\_report.txt} & the checks of the last build
(Section 5) \\
\texttt{references.bib} & the copy of \texttt{src/references.bib} that
the build uses \\
\texttt{README.md}, \texttt{.tex}, \texttt{.pdf} & this description \\
\end{longtable}

\hypertarget{how-the-numbers-are-made}{%
\section{4. How the numbers are made}\label{how-the-numbers-are-made}}

The article prints results of the programs of the folder
\texttt{Revision\_HS\_CJT\_and\_Kohn-Sham}: the Rust programs
\texttt{sf\_part2}, \texttt{sf\_part3} and \texttt{sf\_part4}, which
write the data files of \texttt{data/}, the Python programs
\texttt{sf\_part2.py}, \texttt{sf\_part3.py} and \texttt{sf\_part4.py},
which compute the same results independently with numpy and LAPACK, and
the notebook \texttt{superfluid.nb} in the Wolfram Language, which
repeats the checks of the statements of the article's Secs. III to VI
(except the mass matrix of the gauge bosons and the time-dependent
calculations) with a third implementation. Every check is written to a
report in \texttt{data/} with its deviation and its tolerance.

\texttt{scripts/build\_paper.py} reads the data files and the reports
and computes every number that the four documents state about these
results: the results themselves (the tables, and ranges such as the
Ginzburg numbers), and every bound of the agreement of two computations
(``the three implementations confirm (14) to \(2\times10^{-11}\)''). A
bound is the largest deviation of the checks it summarizes, rounded
\textbf{up} to one significant digit, so that a stated bound is never
smaller than a measured deviation. Typed in the templates are only the
parameters of the model and of the numerical method, mathematical
constants, and the measured values of cosmology, each with its
reference. The values inserted by the last build are listed in
\texttt{generated/numbers.csv}.

The build compiles the article and the Supplemental Material with
pdfLaTeX, BibTeX and REVTeX 4.2, the cover letter and the submission
information with pdfLaTeX, and writes the Markdown versions with pandoc.
It fixes the date that pdfTeX writes into a PDF to the date of the
documents (21 November 2026), so that a rebuild from the same sources
and data gives the same bytes; this was checked by building twice and
comparing the SHA-256 checksums of every PDF, LaTeX and Markdown file.

To build (from the folder \texttt{Revision\_HS\_CJT\_and\_Kohn-Sham},
after the programs have written \texttt{data/}):

\begin{Shaded}
\begin{Highlighting}[]
\BuiltInTok{cd} \StringTok{"/Users/nsh/Developer/github/spacetimes/Revision\_HS\_CJT\_and\_Kohn{-}Sham"}
\ExtensionTok{python3}\NormalTok{ paper\_prd/scripts/build\_paper.py             }\CommentTok{\# the four documents and submission/}
\FunctionTok{bash}\NormalTok{ code/scripts/build\_docs.sh paper\_prd/README.md  }\CommentTok{\# this README as .tex and .pdf}
\end{Highlighting}
\end{Shaded}

\texttt{build\_paper.py} needs Python 3, a TeX installation with REVTeX
4.2 (\texttt{revtex4-2.cls}, \texttt{apsrev4-2.bst}) and the LaTeX
packages \texttt{listings} and \texttt{longtable}, and pandoc. Its exit
status is 1 if a document has an undefined reference or an overfull box
wider than 1 pt, or if the folder \texttt{submission/} fails its check.

\hypertarget{the-checks-of-the-last-build}{%
\section{5. The checks of the last
build}\label{the-checks-of-the-last-build}}

\begin{verbatim}
build_paper.py: the checks of the LaTeX builds (pages; undefined references or citations; overfull boxes wider
than 1 pt).  A document passes when it has no undefined reference and no such box.

higgs_composite_qg: 10 pages; undefined: none; overfull boxes wider than 1 pt: none -> PASS
higgs_composite_qg_SM: 36 pages; undefined: none; overfull boxes wider than 1 pt: none -> PASS
cover_letter: 2 pages; undefined: none; overfull boxes wider than 1 pt: none -> PASS
submission/: the flat source compiles alone to 10 pages (the manuscript has 10); undefined: none -> PASS
submission_info: 6 pages; undefined: none; overfull boxes wider than 1 pt: none -> PASS
\end{verbatim}

The check of \texttt{submission/} compiles a copy of the folder alone in
a temporary folder, with the figures beside the source as the server
keeps them, and compares the number of pages with the article.

\hypertarget{corrections-made-before-the-first-version-was-kept}{%
\section{6. Corrections made before the first version was
kept}\label{corrections-made-before-the-first-version-was-kept}}

Before the first commit, every number typed in the drafts of the
templates was compared with the reports and data files, and every number
that a program computes was replaced by a placeholder. The comparison
found these errors in the drafts, which were not committed (their causes
are recorded in the task record of the folder, entry E20):

\begin{enumerate}
\def\labelenumi{\arabic{enumi}.}
\tightlist
\item
  ``the masses of \(W\) and \(Z\) are at most a fifth of
  \(M^{X,{\rm H}}_{\rm SF}/2\)'' was false for KS and CJT, where the
  ratio is \(0.213\); the text now gives the bound
  \(M_V\le0.11\,M^{X,{\rm H}}_{\rm SF}\) from the data, with \(M_V\) the
  mass of \(W\) or \(Z\), and the tree-level factor
  \((1-4\zeta_V)^{1/2}(1-4\zeta_V+12\zeta_V^2)\),
  \(\zeta_V=M_V^2/(M^{X,{\rm H}}_{\rm SF})^2\), of the decays into
  \(W^+W^-\) and \(ZZ\) (the drafts named only the phase-space factor
  \((1-4\zeta_V)^{1/2}\)). The drafts took the sentence from the
  documents \texttt{superfluid.md} (Section 14.6),
  \texttt{provenance/P4\_remarks.md} (addendum A2) and
  \texttt{superfluid.nb} (Section III.G.8), where it has since been
  corrected (entries E20, E23 and E24 of the task record);
\item
  the ratio of the mass of the Higgs boson to the scale of the breaking
  in the Standard Model was given as 0.509 in the abstract;
  \(125.20/246.22=0.50849\) rounds to 0.508 (0.509 came from rounding
  0.5085 a second time); the abstract then gave 0.5085, and since the
  correction of item 7 it gives no value of the Standard Model;
\item
  ``the programs compute both numbers also from the full static response
  of the normal vacuum \ldots{} to \(10^{-14}\)'': only the Rust program
  computes them so, and the largest deviation is \(3.5\times10^{-14}\);
\item
  bounds smaller than the largest deviation they summarize: ``the
  programs confirm (14) to \(10^{-11}\)'' (the Python check YP3-G2
  deviates by \(1.6\times10^{-11}\)), and, rounded down, ``they agree to
  \(5\times10^{-12}\) (checks P3-G10)'' for the closed form of the slope
  \(\alpha^{X,{\rm GL}}_{\rm SF}\) (the checks P3-G10 deviate by up to
  \(5.1\times10^{-12}\); the sentence now also cites the checks YP3-G10
  and WP3-G10, which deviate by up to \(4.7\times10^{-12}\) and
  \(7.0\times10^{-12}\));
\item
  the build wrote the dash of a range of numbers as
  \texttt{\textbackslash{}text\{-\/-\}}, which pandoc dropped outside
  formulas: the Markdown versions of the drafts read ``Ginzburg numbers
  7.611'' instead of 7.6 to 11; the dash is now written
  \texttt{\textbackslash{}mbox\{-\/-\}}, which LaTeX and pandoc both set
  correctly;
\item
  the Markdown versions of the drafts lacked the title, the byline and
  the abstract, which pandoc does not write for a REVTeX document; the
  build now adds them;
\item
  after the first version had been kept, the author stated (21 November
  2026) that the spinor field, the DiracLike spinor, has electric charge
  0 and no electroweak quantum numbers, and that
  \(SU(2)_{\rm v}\times U(1)^{\rm F}_{\rm LL}\) is a new gauge group,
  not the electroweak group. The first version had called the unbroken
  charge the electric charge and the massless gauge boson the photon,
  and had judged the superfluid against the Higgs boson of the Standard
  Model (the ratio of the mass to the scale of the breaking against
  0.5085, the unitarity bound of the Higgs doublet, the values in GeV on
  the scale \(v=246.22\) GeV, ``charged scalars''). That comparison had
  no basis and is withdrawn: the article, the Supplemental Material and
  the cover letter now answer separately whether the superfluid is the
  Higgs field of its own group and whether it can be that of the
  Standard Model, Table III has no column of the Standard Model and no
  rows in GeV, the vacuum energy is compared with the dark-energy
  density as a condition on the unknown scale, and the massless gauge
  boson is the zeugon (entries C18, C19, E25 to E27 of the task record).
\end{enumerate}

\hypertarget{provenance}{%
\section{7. Provenance}\label{provenance}}

The four documents, the folder \texttt{submission/},
\texttt{generated/numbers.csv} and \texttt{generated/build\_report.txt}
are written by \texttt{scripts/build\_paper.py} from the templates in
\texttt{src/} and the data files and reports of \texttt{data/}; the
Supplemental Material lists the data files, the programs and the
templates with their SHA-256 checksums. The checksums of the templates,
of the build program and of the files to upload, at the time this README
was built:

\begin{itemize}
\tightlist
\item
  \texttt{paper\_prd/scripts/build\_paper.py}:\\
  \texttt{2d493a0239d7ee629a136db0a271b2fdf067b2864962c422019ee289e23aa57b}
\item
  \texttt{paper\_prd/src/higgs\_composite\_qg.template.tex}:\\
  \texttt{d7e42f13b65121f4c6b9edbd47fb6789cee86a2dfe88120391630bcba874db99}
\item
  \texttt{paper\_prd/src/higgs\_composite\_qg\_SM.template.tex}:\\
  \texttt{4fd15b1e7ed260a94f23e0744e8ef63a28de55d4d0f63a0c7840360201c26afd}
\item
  \texttt{paper\_prd/src/cover\_letter.template.tex}:\\
  \texttt{c9ab7ba1a5fdc2f5a7cb8cab08a7470f241b6ed658b0312901c35d00e3afa04f}
\item
  \texttt{paper\_prd/src/submission\_info.template.tex}:\\
  \texttt{a208aec1b66d53011bdccaab5182f48211cdccfc5578a80fcbafccc0654d3773}
\item
  \texttt{paper\_prd/src/references.bib}:\\
  \texttt{3197563f0734ae4cfde3374586d56f4b783bac025ff09e1248c7936cb6c1b1fe}
\item
  \texttt{paper\_prd/submission/higgs\_composite\_qg.pdf}:\\
  \texttt{f8f78ce49b38e0ede09ce492d8038c266ca5925b217b3b2c20449f9479f257e0}
\item
  \texttt{paper\_prd/submission/higgs\_composite\_qg.tex}:\\
  \texttt{be51e0ee904836cbd5a8e33a4247dcb9a456718cd833760755b7142435de93ac}
\item
  \texttt{paper\_prd/submission/higgs\_composite\_qg.bbl}:\\
  \texttt{05cf62568724668752e04505b94c9e766ab1a7046849dbc65adaec0223f1b79d}
\item
  \texttt{paper\_prd/submission/higgs\_composite\_qg\_SM.pdf}:\\
  \texttt{68580df864e4b76359fd0b3f8f0dc61f412021663494ae6a3422932d147c0b8b}
\item
  \texttt{paper\_prd/submission/cover\_letter.pdf}:\\
  \texttt{7a7af727c5d82a82015bdb7a00294e8c9ee5624b352655cc63467232286693e9}
\end{itemize}

\end{document}
